% УВОД, без номер. Ориентировъчен обем: 1 до 2 стр.
% Резюме на дейностите по проекта: защо е поставена задачата, каква е целта,
% какви са задачите, какъв е обхватът и как е устроено изложението.
\section*{УВОД}
\label{sec:introduction}
\addcontentsline{toc}{section}{УВОД}

Достъпът до релационна база данни от програма на C или C++ почти винаги минава през
чужд слой. Приложението извиква клиентска библиотека на доставчика, тя от своя страна
говори с процеса на сървъра по мрежов протокол, а между двете нерядко стои и диспечер
на драйвери като ODBC. Всеки слой има своя цена: собствен модел на паметта, собствено
разбиране за типовете, собствено копие на резултатните редове и собствен, обикновено
синхронен, модел на изчакване. Приложението плаща тази цена, без да я вижда, защото
слоят е даденост.

Настоящият курсов проект разглежда алтернативата. Мрежовите протоколи на PostgreSQL и
на MySQL са публично специфицирани~\cite{pg_protocol, mysql_protocol}, което означава,
че клиентът може да бъде написан срещу самата спецификация, а не срещу чужда
библиотека. Този избор премахва междинните слоеве и прави видими решения, които иначе
са скрити: къде се копира буфер, кога се сериализира стойност, колко обръщания до
сървъра струва едно подготвено заявление. Същият избор носи и очевидни рискове, които
проектът е длъжен да отчете честно, а не да премълчи.

\textbf{Целта} на проекта е да се проектира и реализира клиентска библиотека на C++20,
която реализира пряко мрежовите протоколи на PostgreSQL и MySQL, и да се измери
поведението ѝ спрямо доставчиковите клиентски библиотеки върху едно и също натоварване.

За постигането на целта са формулирани следните \textbf{задачи}:

\begin{enumerate}[topsep=0pt,itemsep=0pt,leftmargin=*]
\item преглед на стандарта SQL и на изградените върху него технологии за достъп,
      както и на спецификациите на двата разглеждани мрежови протокола;
\item анализ на възможните решения, слоест диспечер на драйвери, доставчикова
      клиентска библиотека или пряка реализация на протокола, и обосновка на избора;
\item проектиране на слоеста архитектура с ясно отделени транспорт, протоколни
      кодеци, декодер на редове, пул от връзки и слой за типово изображение;
\item програмна реализация на архитектурата за двата протокола;
\item изграждане на методика за измерване и на възпроизводима среда за нея;
\item експериментално сравнение с \code{libpq} и \code{libmysqlclient} и анализ на
      получените резултати.
\end{enumerate}

\textbf{Обхватът} на проекта е клиентската страна на комуникацията: установяване на
връзка, удостоверяване, изпълнение на прости и на подготвени заявления, декодиране на
резултатни редове и управление на пула от връзки. Извън обхвата остават
администрирането на сървъра, репликацията, разпределените транзакции и всяка форма на
съхранение от страна на клиента.

\textbf{Структура на изложението.} Раздел~\ref{sec:literature} представя предметната
област и прегледа на литературата. Раздел~\ref{sec:alternatives} излага разгледаните
алтернативи и обосновава избора. Раздел~\ref{sec:design} описва проектирането, а
Раздел~\ref{sec:implementation} програмната реализация. Раздел~\ref{sec:method}
дефинира методиката за измерване, а Раздел~\ref{sec:results} представя резултатите и
техния анализ. Заключението обобщава постигнатото и посочва посоките за продължение.
